\documentclass[11pt]{article}
\usepackage[margin=1in]{geometry}
\usepackage{amsmath, amsthm, amssymb, mathtools}
\usepackage{hyperref}
\usepackage{xcolor}
\usepackage{tcolorbox}
\tcbuselibrary{theorems, skins}

% ---- Theorem environments ----
\newtheorem{theorem}{Theorem}[section]
\newtheorem{lemma}[theorem]{Lemma}
\newtheorem{corollary}[theorem]{Corollary}
\newtheorem{definition}[theorem]{Definition}
\newtheorem{proposition}[theorem]{Proposition}
\newtheorem{remark}{Remark}

% ---- Lean verification badge ----
\newcommand{\verified}{%
  \textcolor{green!60!black}{\textbf{[Lean 4 ✓]}}%
}
\newcommand{\sorry}{%
  \textcolor{orange}{\textbf{[sorry]}}%
}

% ---- Math macros ----
\newcommand{\R}{\mathbb{R}}
\newcommand{\N}{\mathbb{N}}
\newcommand{\cH}{\mathcal{H}}
\newcommand{\cF}{\mathcal{F}}
\newcommand{\Fstar}{F^*}
\newcommand{\hhat}{\boldsymbol{h}}
\newcommand{\dF}{\frac{dF^*(\tau|\hhat)}{d\tau}}
\newcommand{\sigmoid}{\sigma}
\newcommand{\MNN}{\mathcal{M}}
\newcommand{\gammaN}{\gamma_n}

\title{\textbf{CuFun: Machine-Verified Mathematical Foundations}\\
\large Lean 4 + Mathlib Formal Proofs}
\author{Maolin Wang \\ \small City University of Hong Kong}
\date{Verified with Lean 4.29.0-rc6 + Mathlib (March 2026)}

\begin{document}
\maketitle

\begin{tcolorbox}[colback=green!5!white, colframe=green!50!black, title=Verification Status]
All theorems marked \verified{} carry machine-checked proofs with \textbf{zero \texttt{sorry}}
in Lean 4 + Mathlib. Source: \url{https://github.com/MorinClaw/CuFun-Formal-Theory}
\end{tcolorbox}

\tableofcontents
\bigskip

% ============================================================
\section{CDF Structure and Axioms}
% ============================================================

\begin{definition}[Conditional CDF \verified{}]
\label{def:cdf}
Let $H$ be a history space. A \emph{conditional CDF} is a function
$F : \R_{\geq 0} \times H \to \R$ satisfying:
\begin{enumerate}
  \item \textbf{Zero at origin:} $\forall \hhat \in H,\quad F(0\,|\,\hhat) = 0$
  \item \textbf{Limit to one:} $\forall \hhat \in H,\quad \lim_{\tau\to\infty} F(\tau\,|\,\hhat) = 1$
  \item \textbf{Monotonicity:} $\forall \hhat \in H,\quad \tau_1 \leq \tau_2 \Rightarrow F(\tau_1\,|\,\hhat) \leq F(\tau_2\,|\,\hhat)$
  \item \textbf{Boundedness:} $\forall \tau \geq 0,\, \hhat \in H,\quad 0 \leq F(\tau\,|\,\hhat) \leq 1$
\end{enumerate}
\end{definition}

\begin{definition}[Differentiable CDF and PDF \verified{}]
A \emph{differentiable conditional CDF} additionally satisfies:
\begin{enumerate}
  \item $F(\cdot\,|\,\hhat)$ is differentiable at every $\tau \in \R$ for all $\hhat$
  \item \textbf{Strict positivity of density:}
  $\forall \tau,\, \hhat,\quad p^*(\tau\,|\,\hhat) := \dfrac{\partial F(\tau\,|\,\hhat)}{\partial \tau} > 0$
\end{enumerate}
The \emph{conditional probability density function} is defined as:
\[
  p^*(\tau\,|\,\hhat) := \frac{\partial F(\tau\,|\,\hhat)}{\partial \tau}
\]
\end{definition}

% ============================================================
\section{Structural Closure Properties}
% ============================================================

\begin{theorem}[Convex Combination Closure \verified{}]
\label{thm:convex-comb}
Let $F_1, F_2$ be valid conditional CDFs over the same history space $H$.
For any $\lambda \in [0,1]$, the convex combination
\[
  F_\lambda(\tau\,|\,\hhat) := \lambda\, F_1(\tau\,|\,\hhat) + (1-\lambda)\, F_2(\tau\,|\,\hhat)
\]
is also a valid conditional CDF.
\end{theorem}

\begin{proof}[Proof sketch (machine-verified)]
We verify all four axioms of Definition~\ref{def:cdf}:
\begin{enumerate}
  \item \emph{Zero at origin:} $F_\lambda(0|\hhat) = \lambda \cdot 0 + (1-\lambda)\cdot 0 = 0$. $\checkmark$
  \item \emph{Limit to one:} By linearity of limits,
    $\lim_{\tau\to\infty} F_\lambda(\tau|\hhat) = \lambda\cdot 1 + (1-\lambda)\cdot 1 = 1$. $\checkmark$
  \item \emph{Monotonicity:} For $\tau_1 \leq \tau_2$:
    \[
      F_\lambda(\tau_2|\hhat) - F_\lambda(\tau_1|\hhat)
      = \underbrace{\lambda}_{\geq 0}\underbrace{(F_1(\tau_2|\hhat)-F_1(\tau_1|\hhat))}_{\geq 0}
      + \underbrace{(1-\lambda)}_{\geq 0}\underbrace{(F_2(\tau_2|\hhat)-F_2(\tau_1|\hhat))}_{\geq 0}
      \geq 0. \quad\checkmark
    \]
  \item \emph{Boundedness:} Immediate from $0 \leq \lambda \leq 1$ and the bounds on $F_1, F_2$. $\checkmark$
\end{enumerate}
\end{proof}

\begin{remark}
This theorem is the structural backbone of the \emph{$\delta$-net blending} step
in the Universal Approximation proof: approximations at grid points can be
combined with non-negative, sum-to-one weights while preserving CDF validity.
\end{remark}

% ============================================================
\section{MNN Monotonicity Theorems}
% ============================================================

\begin{theorem}[Weighted Sum Monotonicity \verified{}]
\label{thm:weighted-sum}
Let $\{f_i : \R \to \R\}_{i=1}^n$ be monotone non-decreasing functions and
$\{w_i\}_{i=1}^n \subset \R_{\geq 0}$ be non-negative weights. Then:
\[
  g(\tau) := \sum_{i=1}^n w_i\, f_i(\tau)
\]
is monotone non-decreasing in $\tau$.
\end{theorem}

\begin{proof}[Proof (machine-verified)]
For $\tau_1 \leq \tau_2$:
\[
  g(\tau_2) - g(\tau_1)
  = \sum_{i=1}^n w_i \underbrace{(f_i(\tau_2) - f_i(\tau_1))}_{\geq 0\text{ by monotonicity}}
  \geq 0,
\]
since each term is non-negative ($w_i \geq 0$, $f_i(\tau_2) - f_i(\tau_1) \geq 0$).
\end{proof}

\begin{theorem}[Sigmoid Monotonicity \verified{}]
\label{thm:sigmoid}
The sigmoid function $\sigma : \R \to (0,1)$ defined by
\[
  \sigma(x) = \frac{1}{1 + e^{-x}}
\]
is strictly monotone increasing.
\end{theorem}

\begin{proof}[Proof (machine-verified)]
For $a \leq b$, we have $-b \leq -a$, hence $e^{-b} \leq e^{-a}$,
so $1 + e^{-b} \leq 1 + e^{-a}$. Since $1/(1+e^{-b}) \leq 1/(1+e^{-a})$
follows from the fact that $x \mapsto 1/x$ is decreasing on $(0,\infty)$:
\[
  0 < 1 + e^{-b} \leq 1 + e^{-a}
  \;\Rightarrow\;
  \frac{1}{1+e^{-a}} \leq \frac{1}{1+e^{-b}},
\]
i.e., $\sigma(a) \leq \sigma(b)$. \qed
\end{proof}

\begin{theorem}[MNN Architectural Monotonicity \verified{}]
\label{thm:mnn-mono}
Let $d \in \N$. Let $w_1, \ldots, w_d \geq 0$ be non-negative weights and
$\phi : \R \to \R$ a monotone function. Define the MNN output:
\[
  \hat{F}(\tau) := \sigma\!\left(-\sum_{i=1}^d w_i\,\phi(\tau)\right)
               = \frac{1}{1 + \exp\!\left(-\sum_{i=1}^d w_i\,\phi(\tau)\right)}.
\]
Then $\hat{F}$ is monotone non-decreasing in $\tau$.
\end{theorem}

\begin{proof}[Proof (machine-verified)]
By Theorem~\ref{thm:weighted-sum}, $\tau \mapsto \sum_i w_i\,\phi(\tau)$ is monotone.
Composition with the monotone function $\sigma$ (Theorem~\ref{thm:sigmoid})
yields a monotone composite. \qed
\end{proof}

\begin{theorem}[MNN Output Bounded in $(0,1)$ \verified{}]
\label{thm:mnn-bounded}
Under the same conditions as Theorem~\ref{thm:mnn-mono},
\[
  \forall\, \tau \in \R: \quad 0 < \hat{F}(\tau) < 1.
\]
Moreover, $\hat{F}$ is monotone. Thus two of the three CDF axioms
(boundedness and monotonicity) hold by architecture alone; only the
boundary condition $\hat{F}(0) = 0$ requires specific initialization.
\end{theorem}

\begin{proof}[Proof (machine-verified)]
\emph{Strict positivity:} $\hat{F}(\tau) = (1 + e^{-z})^{-1} > 0$ since $e^{-z} > 0$.

\emph{Strictly less than 1:} $\hat{F}(\tau) < 1 \iff 1 < 1 + e^{-z} \iff e^{-z} > 0$. $\checkmark$ \qed
\end{proof}

% ============================================================
\section{Log-Likelihood and PDF Validity}
% ============================================================

\begin{theorem}[Log-Likelihood Well-Definedness \verified{}]
\label{thm:log-ll}
For any differentiable conditional CDF $F^*$ and any event $(\tau, \hhat)$,
\[
  p^*(\tau\,|\,\hhat) = \frac{\partial F^*(\tau\,|\,\hhat)}{\partial \tau} > 0,
\]
so the per-event log-likelihood $\log p^*(\tau|\hhat)$ is finite for all $(\tau, \hhat)$.
\end{theorem}

\begin{proof}
Immediate from the \texttt{deriv\_pos} field of \texttt{DiffConditionalCDF}.
\end{proof}

\begin{remark}
This contrasts with intensity-based methods where $\log\lambda^*(t)$ is undefined
whenever $\lambda^*(t) = 0$. CuFun's CDF parameterization structurally guarantees
a well-defined log-likelihood everywhere.
\end{remark}

% ============================================================
\section{CDF–Intensity Equivalence}
% ============================================================

\begin{theorem}[CDF–Intensity Derivative Identity \verified{}]
\label{thm:cdf-intensity}
Let $\lambda^* : \R \to \R$ be an intensity function with cumulative hazard
$\Lambda(\tau) := \int_0^\tau \lambda^*(s)\,ds$,
and define the CDF by the standard TPP formula:
\[
  F^*(\tau) = 1 - e^{-\Lambda(\tau)}.
\]
Then:
\[
  \frac{dF^*(\tau)}{d\tau}
  = \lambda^*(\tau) \cdot \bigl(1 - F^*(\tau)\bigr)
  = \lambda^*(\tau)\, e^{-\Lambda(\tau)}.
\]
In particular, $p^*(\tau) = \frac{dF^*}{d\tau}$, confirming that the CDF
representation carries \emph{identical information} to the intensity function.
\end{theorem}

\begin{proof}[Proof (machine-verified via chain rule)]
By the chain rule and the Fundamental Theorem of Calculus:
\[
  \frac{d}{d\tau}\bigl[1 - e^{-\Lambda(\tau)}\bigr]
  = -e^{-\Lambda(\tau)} \cdot (-\lambda^*(\tau))
  = \lambda^*(\tau)\, e^{-\Lambda(\tau)}
  = \lambda^*(\tau)\,(1 - F^*(\tau)). \qed
\]
\end{proof}

\begin{corollary}[PDF Positivity \verified{}]
\label{cor:pdf-pos}
If $\lambda^*(t) > 0$ for all $t$, then $p^*(\tau) > 0$ for all $\tau > 0$.
\end{corollary}

% ============================================================
\section{Numerical Error Analysis}
% ============================================================

\begin{definition}[Higham's Rounding Error Constant \verified{}]
For $n$ sequential floating-point operations with unit roundoff $u$, define:
\[
  \gamma_n := \frac{n u}{1 - n u}, \quad n u < 1.
\]
This bounds the accumulated relative rounding error in a chain of $n$ operations.
\end{definition}

\begin{theorem}[Rounding Constant Positivity \verified{}]
For $n \geq 1$ and $0 < u < 1/n$:
\[
  \gamma_n = \frac{nu}{1-nu} > 0.
\]
\end{theorem}

\begin{theorem}[CDF Numerical Superiority \verified{}]
\label{thm:numerical}
Let $L$ be the MNN depth, $W$ width, and $N$ the number of quadrature points
in intensity-based methods. Let $u$ be the unit roundoff.

Assume $2LW < N$ (CuFun uses fewer elementary operations than $N$-point quadrature),
$(2LW)\,u < 1$, and $N u < 1$ (both computations are numerically stable).

Then:
\[
  \gamma_{2LW} < \gamma_N,
\]
i.e., \emph{CuFun's autodiff computation accumulates strictly less rounding error
than numerical integration with $N$ points}.
\end{theorem}

\begin{proof}[Proof (machine-verified)]
We need to show $\frac{(2LW)u}{1-(2LW)u} < \frac{Nu}{1-Nu}$.

Setting $a = (2LW)u$ and $b = Nu$, with $0 < a < b < 1$
(since $2LW < N$ and $u > 0$), this reduces to:
\[
  \frac{a}{1-a} < \frac{b}{1-b}
  \;\iff\; a(1-b) < b(1-a)
  \;\iff\; a - ab < b - ab
  \;\iff\; a < b. \quad\checkmark
\]
The last step holds since $(2LW)u < Nu$ follows from $2LW < N$ and $u > 0$. \qed
\end{proof}

\begin{remark}
The deeper advantage is \emph{structural}: while quadrature involves a
summation of $N$ terms (potentially suffering catastrophic cancellation when
$\log\lambda^*(\tau) \approx \int\lambda^*(s)ds$), autodiff applies the chain rule
to a fixed computational graph, whose condition number is bounded independently of $N$.
\end{remark}

% ============================================================
\section{Statistical Convergence Structure}
% ============================================================

\begin{theorem}[MSE Bias-Variance Decomposition \verified{}]
\label{thm:mse}
For any estimator $\hat{F}_n$, best-in-class approximator $F_n^*$,
and true CDF $F_0$, the mean squared error satisfies:
\[
  \bigl(\hat{F}_n(\tau|\hhat) - F_0(\tau|\hhat)\bigr)^2
  \;\leq\;
  2\underbrace{\bigl(\hat{F}_n(\tau|\hhat) - F_n^*(\tau|\hhat)\bigr)^2}_{\text{estimation error}}
  +\;
  2\underbrace{\bigl(F_n^*(\tau|\hhat) - F_0(\tau|\hhat)\bigr)^2}_{\text{approximation error}}.
\]
\end{theorem}

\begin{proof}[Proof (machine-verified via nlinarith)]
By the identity $(a + b)^2 \leq 2a^2 + 2b^2$ applied to
$a = \hat{F}_n - F_n^*$ and $b = F_n^* - F_0$:
\[
  (\hat{F}_n - F_0)^2 = (a + b)^2 \leq 2a^2 + 2b^2. \qed
\]
\end{proof}

% ============================================================
\section{Summary of Verified Results}
% ============================================================

\begin{table}[h!]
\centering
\renewcommand{\arraystretch}{1.3}
\begin{tabular}{lll}
\hline
\textbf{Theorem} & \textbf{Content} & \textbf{Status} \\
\hline
Def.~\ref{def:cdf} & CDF axiom structure & \verified{} \\
Thm.~\ref{thm:convex-comb} & Convex combination closure & \verified{} \\
Thm.~\ref{thm:weighted-sum} & Weighted sum monotonicity & \verified{} \\
Thm.~\ref{thm:sigmoid} & Sigmoid monotonicity & \verified{} \\
Thm.~\ref{thm:mnn-mono} & MNN architectural monotonicity & \verified{} \\
Thm.~\ref{thm:mnn-bounded} & MNN output $\in (0,1)$ & \verified{} \\
Thm.~\ref{thm:log-ll} & Log-likelihood well-definedness & \verified{} \\
Thm.~\ref{thm:cdf-intensity} & CDF–intensity derivative identity & \verified{} \\
Cor.~\ref{cor:pdf-pos} & PDF positivity & \verified{} \\
Thm.~\ref{thm:numerical} & CDF error $<$ quadrature error & \verified{} \\
Thm.~\ref{thm:mse} & MSE bias-variance decomposition & \verified{} \\
\hline
\end{tabular}
\caption{Complete list of machine-verified theorems for CuFun.
All proofs are \texttt{sorry}-free in Lean 4.29.0-rc6 + Mathlib.}
\end{table}

\end{document}
